\documentclass[10pt,a4paper,twocolumn]{article}

\usepackage[T2A]{fontenc}
\usepackage[utf8]{inputenc}
\usepackage[russian,english]{babel}
\usepackage{amsmath}
\usepackage{array}
\usepackage{booktabs}
\usepackage{graphicx}
\usepackage{geometry}
\usepackage{hyperref}
\usepackage{microtype}
\usepackage[numbers,sort&compress]{natbib}
\usepackage{xcolor}

\geometry{
  top=17mm,
  bottom=19mm,
  left=17mm,
  right=17mm,
  columnsep=7mm
}
\AtBeginDocument{\selectlanguage{russian}\fontencoding{T2A}\selectfont}
\hypersetup{
  colorlinks=true,
  linkcolor=blue!45!black,
  citecolor=blue!45!black,
  urlcolor=blue!45!black
}

\newcommand{\mbar}{\overline{m}}
\newcommand{\Mbar}{\overline{M}}
\newcommand{\paperfigure}[2]{%
  \IfFileExists{#1}{%
    \includegraphics[width=#2]{#1}%
  }{%
    \fbox{\parbox[c][42mm][c]{0.90\linewidth}{\centering
      Рисунок ещё не построен.\\
      Запустите соответствующую ячейку в
      \texttt{code/sbf-2-graph.ipynb}.\\[1mm]
      \texttt{\detokenize{#1}}}}%
  }%
}

\title{Прямая калибровка флуктуаций поверхностной яркости
JWST/NIRCam F150W по выборке TRGB программы GO-3055}
\author{И.\,А.~Буканов\\[-1mm]
\small Московский государственный университет имени М.\,В.~Ломоносова,\\[-1mm]
\small факультет космических исследований, Москва, Россия\\[1mm]
\small Состав соавторов уточняется}
\date{Черновик от 1 августа 2026 года}

\begin{document}
\maketitle

\begin{abstract}
Представлена прямая эмпирическая калибровка флуктуационных звёздных величин
(SBF), измеренных по изображениям JWST/NIRCam в фильтре F150W. Выборка
содержит 14 галактик ранних морфологических типов из программы GO-3055,
наблюдавшихся для проекта шкалы расстояний TRGB--SBF. Геометрия измерения
следует Jensen et al. (2025): SBF и цвет определяются в двух концентрических
круговых кольцах $8.2$--$16.4$ и $16.4$--$32.8$ угл.\ сек, после чего два
измерения объединяются с весами, обратными дисперсиям. Цвет F090W--F150W
измеряется по тем же пикселям, что и SBF-сигнал. Абсолютные флуктуационные
величины привязаны к пересмотренным цветозависимым расстояниям JWST TRGB из
Chazov et al. (2026). Цвет, SBF-величина и апертурная светимость исправлены
за поглощение в Млечном Пути, а ковариация этой поправки включена в
правдоподобие. Для полной выборки получено
\[
\begin{aligned}
 \Mbar_{150,0}={}&(-3.187^{+0.026}_{-0.020})\\
 &+(0.523^{+0.484}_{-0.577})
 [C_{090-150,0}-0.5675],
\end{aligned}
\]
при внутреннем разбросе $0.068$ зв.\ вел. Перекрёстная проверка с исключением
по одной галактике даёт среднее отличие от TRGB $+0.006$ зв.\ вел. и
RMS $0.099$ зв.\ вел., то есть около $4.7\%$ по расстоянию. Модель без
цветового члена умеренно предпочтительнее
($\Delta{\rm AICc}=-2.45$ относительно линейной модели), поэтому имеющаяся
выборка не даёт однозначного обнаружения цветовой зависимости. Внутреннее и
внешнее кольца показывают различные тренды; в частности, для NGC~3379
получена разница $0.168$ зв.\ вел. при исключительно однородном остаточном
изображении. Обсуждаются ограничения текущего измерения и условия переноса
TRGB-нуль-пункта на будущие выборки JWST SBF.
\end{abstract}

\noindent\textbf{Ключевые слова:} шкала расстояний; галактики; флуктуации
поверхностной яркости; JWST; NIRCam; TRGB.

\section{Введение}

Флуктуации поверхностной яркости возникают из-за пуассоновских изменений
числа и светимости неразрешённых звёзд в отдельных элементах разрешения
\citep{TonrySchneider1988}. В ближнем инфракрасном диапазоне SBF-сигнал старых
звёздных популяций достаточно велик, чтобы измерять расстояния до галактик
ранних типов далеко за пределами области, где доступны индивидуальные
звёзды. Однако абсолютная амплитуда SBF зависит от функции светимости
эволюционировавших звёзд и должна калиброваться по наблюдаемому индикатору
звёздного населения, обычно по интегральному цвету.

Проект TRGB--SBF строит шкалу расстояний населения II, в которой геометрически
калиброванный нуль-пункт вершины ветви красных гигантов (TRGB) переносится на
метод SBF. Jensen et al. \citep[далее Paper III;][]{Jensen2025} использовали
расстояния JWST TRGB для галактик GO-3055, чтобы уточнить существующую
калибровку HST/WFC3-IR F110W SBF. При использовании средних расстояний
скоплений Дева и Форнакс они получили
\[
 \Mbar_{110}=(1.86\pm0.16)[(g-z)_{\rm PS}-1.30]
             -(2.760\pm0.016).
\]

В настоящей работе и флуктуационный сигнал, и индикатор звёздного населения
измеряются непосредственно по изображениям JWST GO-3055. Фильтр F150W
используется как сигнальный, а F090W--F150W --- как цвет. В качестве опорных
расстояний приняты значения для тонкого среза из Chazov et al.
\citep[далее Paper IV;][]{Chazov2026}. Их калибровка F090W TRGB опирается на
мазерное расстояние до NGC~4258. Для
$(F090W-F150W)_0<1.65$ зв.\ вел. они получили
$M_{\rm F090W}^{\rm TRGB}=-4.398\pm0.017$ зв.\ вел.; при более красных цветах
используется полиномиальная поправка. Новые расстояния тонкого среза в среднем
на $0.028\pm0.008$ зв.\ вел. меньше прежних литературных значений для тех же
изображений.

Цели работы: воспроизвести геометрию измерений Paper III, получить прямую
калибровку JWST F150W SBF, проверить восстановленные расстояния относительно
TRGB и выяснить, объясняют ли радиальные изменения цвета и светимость
галактики часть остаточного разброса.

\section{Данные}

\subsection{Выборка GO-3055}

Выборка включает 14 галактик: три объекта скопления Форнакс, девять объектов
области Девы и две галактики близких групп. Для каждой цели имеются
изображения NIRCam F090W и F150W. Величина SBF измеряется в F150W, а
$C_{090-150}=F090W-F150W$ используется как индикатор звёздного населения.
Принятые модули расстояния TRGB для тонкого среза из Paper IV привязаны к
мазерному модулю NGC~4258, равному $29.397$ зв.\ вел.

\begin{table*}
\centering
\caption{Измерения для калибровки F150W SBF по выборке GO-3055.
Предпоследний столбец содержит остаток SBF-расстояния при исключении данной
галактики из калибровки относительно принятого модуля TRGB из Paper IV. Знак
HQ IV показывает флаг high-quality из таблицы~2 Paper IV. Указанная
$\sigma_{\overline M}$ содержит индивидуальные, но не общие нуль-пункты.}
\label{tab:measurements}
\resizebox{\textwidth}{!}{%
\begin{tabular}{llrrrrrrrc}
\toprule
Галактика & Окружение & $C_{090-150,0}$ & $\overline m_{150,0}$ &
$\overline M_{150,0}$ & $\sigma_{\overline M}$ & $|\Delta\overline m|_{\text{кол.}}$ &
$\mu_{\rm TRGB}$ & $\Delta\mu_{\rm LOO}$ & HQ IV \\
 & & (зв.\ вел.) & (зв.\ вел.) & (зв.\ вел.) &
(зв.\ вел.) & (зв.\ вел.) & (зв.\ вел.) & (зв.\ вел.) & \\
\midrule
NGC 1380 & Форнакс & 0.510 & 28.111 & -3.297 & 0.034 & 0.046 & 31.408 & -0.115 & $+$ \\
NGC 1399 & Форнакс & 0.570 & 28.180 & -3.318 & 0.034 & 0.019 & 31.498 & -0.143 & $-$ \\
NGC 1404 & Форнакс & 0.559 & 28.206 & -3.160 & 0.028 & 0.026 & 31.366 & +0.035 & $+$ \\
NGC 1549 & другая группа & 0.533 & 27.986 & -3.245 & 0.035 & 0.046 & 31.231 & -0.049 & $-$ \\
NGC 3379 & другая группа & 0.557 & 27.042 & -3.036 & 0.089 & 0.168 & 30.078 & +0.164 & $+$ \\
NGC 4374 & Дева & 0.566 & 27.965 & -3.147 & 0.035 & 0.043 & 31.112 & +0.044 & $+$ \\
NGC 4406 & Дева & 0.616 & 27.863 & -3.233 & 0.041 & 0.062 & 31.096 & -0.085 & $+$ \\
NGC 4472 & Дева & 0.625 & 27.763 & -3.253 & 0.035 & 0.039 & 31.016 & -0.122 & $-$ \\
NGC 4486 & Дева & 0.652 & 27.878 & -3.168 & 0.062 & 0.069 & 31.046 & -0.035 & $-$ \\
NGC 4552 & Дева & 0.569 & 27.876 & -3.047 & 0.051 & 0.086 & 30.923 & +0.150 & $+$ \\
NGC 4621 & Дева & 0.565 & 27.717 & -3.104 & 0.044 & 0.072 & 30.821 & +0.092 & $+$ \\
NGC 4636 & Дева & 0.604 & 27.926 & -3.144 & 0.031 & 0.022 & 31.070 & +0.026 & $+$ \\
NGC 4649 & Дева & 0.623 & 27.961 & -3.048 & 0.064 & 0.119 & 31.009 & +0.127 & $-$ \\
NGC 4697 & Дева & 0.512 & 27.109 & -3.215 & 0.064 & 0.114 & 30.324 & +0.000 & $+$ \\
\bottomrule
\end{tabular}
}
\end{table*}

\subsection{Связь с работами Paper III и Paper IV}

В таблице~\ref{tab:method-comparison} разделены различия, вызванные фильтрами
и инструментами, и собственно методические решения. Главный ещё не закрытый
тест --- повторить калибровку с точным использованием средних расстояний
скоплений из Paper III. В текущем варианте применяются индивидуальные
расстояния Paper IV.

\begin{table*}
\centering
\caption{Связь настоящего анализа с Paper III и Paper IV.}
\label{tab:method-comparison}
\resizebox{\textwidth}{!}{%
\begin{tabular}{
>{\raggedright\arraybackslash}p{0.17\textwidth}
>{\raggedright\arraybackslash}p{0.25\textwidth}
>{\raggedright\arraybackslash}p{0.25\textwidth}
>{\raggedright\arraybackslash}p{0.25\textwidth}}
\toprule
Элемент & Paper III: Jensen et al. (2025) &
Paper IV: Chazov et al. (2026) & Настоящая работа \\
\midrule
Основная величина &
HST/WFC3-IR F110W SBF &
JWST/NIRCam F090W TRGB &
JWST/NIRCam F150W SBF \\
Индикатор населения &
DECam или Pan-STARRS $g-z$ &
цвет отдельных звёзд F090W--F150W &
интегральный F090W--F150W в SBF-кольцах \\
Опорное расстояние &
JWST TRGB; преимущественно средние Девы и Форнакса &
мазер NGC~4258; индивидуальные TRGB для тонкого и толстого срезов &
индивидуальные TRGB тонкого среза Paper IV \\
Радиальные области &
$8.2$--$16.4$ и $16.4$--$32.8$ угл.\ сек; взвешенное среднее &
внешние поля разрешённых звёзд &
те же два кольца; среднее и покольцевая диагностика \\
Аппроксимация цвета &
байесовская, с погрешностями по обеим осям &
кусочный цветовой закон TRGB с границей $1.65$ зв.\ вел. &
правдоподобие с внутренним разбросом, ошибкой цвета и её ковариацией с SBF \\
Размер калибровки &
14 объектов; восемь прямых пересечений с TRGB &
16 целей по расстояниям, 12 основных для цветового закона &
14 прямых пар SBF--TRGB \\
\bottomrule
\end{tabular}
}
\end{table*}

\begin{table}
\centering
\caption{Параметры TRGB из Paper IV, используемые в калибровке SBF.}
\label{tab:paper4-anchor}
\resizebox{\columnwidth}{!}{%
\begin{tabular}{lc}
\toprule
Величина & Значение Paper IV \\
\midrule
Синий нуль-пункт TRGB & $-4.398\pm0.017$ \\
Граница по цвету & $1.65$ \\
Коэффициенты $(A,B,C,D)$ &
$(-4.398,\,0.431,\,0.712,\,-0.293)$ \\
Тонкий минус прежний модуль & $-0.028\pm0.008$ \\
Толстый минус прежний модуль & $-0.037\pm0.009$ \\
Среднее тонкий--толстый & $0.007$ \\
Общая систематика нуль-пункта & $0.047$ \\
Принятый вариант & тонкий срез \\
\bottomrule
\end{tabular}
}
\end{table}

\section{Измерение SBF}

\subsection{Фон, маски и модель галактики}

По области вне рабочих колец оценивается и вычитается скалярный фон F150W.
Первичная маска исключает невалидные пиксели, щели детекторов, объекты переднего
плана и компактные источники. По незамаскированному свету галактики строятся
изофоты, которые интерполируются в гладкую полноразмерную модель. Щели между
детекторами заполняются только для устойчивого построения изофот; их пиксели
не участвуют в измерении SBF.

Модель вычитается из изображения после вычитания фона. Компактные источники
повторно детектируются по остаткам и маскируются. Рабочий остаток нормируется
на квадратный корень из гладкой модели. Симметричное отсечение крайних
значений на уровне $3.5\sigma$ сохранено как сознательное отличие от
исторического пайплайна. Его влияние на нуль-пункт должно быть представлено
отдельным тестом чувствительности.

\subsection{Спектр мощности и неразрешённые источники}

В каждом кольце азимутально усреднённый спектр мощности остатка
аппроксимируется на интервале $0.04\leq k\leq0.25$:
\begin{equation}
 P(k)=P_0E(k)+P_1.
\end{equation}
Здесь $E(k)$ --- ожидаемая форма SBF-сигнала, полученная из PSF с учётом
маски и окна кольца; $P_0$ --- амплитуда компоненты этой формы; $P_1$ ---
постоянная компонента белого шума. Мощность, приписываемая звёздной популяции,
равна
\begin{equation}
 P_{\rm SBF}=P_0-P_r,
\end{equation}
где $P_r$ оценивает вклад неразрешённых шаровых скоплений и фоновых галактик
слабее порога маскирования. Кажущаяся флуктуационная звёздная величина
вычисляется из $P_{\rm SBF}$, фотометрического нуль-пункта и среднего уровня
модели галактики в кольце.

Во всех текущих покольцевых фитах $P_1<0$. Прямой контроль четырёх угловых
областей исходных i2d показывает не белый шум: медианный наклон
$\log N(k)$ в диапазоне $0.04<k<0.25$ равен $-1.45$, а отношение мощности
в высокочастотной и низкочастотной частях диапазона --- $0.18$. Поэтому
постоянный $P_1$ является приближением. Диагностика подтверждает необходимость
будущего фита с явным шаблоном $N(k)$, но сама по себе не задаёт его форму в
рабочих кольцах галактики.

\subsection{Цвет и погрешности}

Изображение F090W согласуется по PSF и системе координат с F150W. Цвет
измеряется по тем же рабочим пикселям и в тех же кольцах, где определяется
SBF. Это принципиально: цвет центральной области галактики не обязан
соответствовать звёздному населению SBF-колец.

Для каждой галактики из таблицы~2 Paper IV взято $A_{090}$ и вычислены
$E(B-V)=A_{090}/1.4156$ и $A_{150}=0.6021E(B-V)$. Используются
\begin{align}
 C_{090-150,0}&=C_{\rm obs}-(A_{090}-A_{150}),\\
 \mbar_{150,0}&=\mbar_{150,\rm obs}-A_{150}.
\end{align}
Поскольку модуль TRGB и наш цвет используют ту же поправку, в правдоподобии
сохранена ковариация
$\operatorname{cov}(C_0,\Mbar_{150,0})=-\sigma^2(A_{090}-A_{150})$.

Индивидуальная погрешность включает ошибку подгонки, устойчивость к
$k_{\min}$, различие колец, PSF-ансамбль, рабочую неопределённость $P_r$,
линеаризованный вклад разброса независимых оценок фона и поглощение.
PSF не суммируется повторно: он уже входит в ошибку $P_0$. Общие систематики
вынесены отдельно: $0.047$ зв.\ вел. для нуль-пункта TRGB/NGC~4258 и
$0.0108$ зв.\ вел. для принятой 1\%-ной абсолютной калибровки NIRCam;
их квадратичная сумма равна $0.048$ зв.\ вел.

\section{Результаты}

\subsection{Калибровка F150W}

Для 14 индивидуальных расстояний Paper IV при медианном исправленном цвете
$C_0=0.5675$ получена калибровка
\begin{equation}
 \Mbar_{150,0}=(-3.1869)
 +(0.5232)(C_{090-150,0}-0.5675).
 \label{eq:f150-calibration}
\end{equation}
Внутренний разброс равен $0.0676$ зв.\ вел. Бутстрэп-интервалы 16--84
процентилей составляют $[-3.206,-3.161]$ для нуль-пункта и
$[-0.054,1.007]$ для наклона. Широкий интервал наклона обусловлен узким
диапазоном цветов и малой выборкой.

\begin{figure*}
\centering
\paperfigure{../../runs/sbf2_go3055/analysis/figures/go3055_f150w_sbf_calibration.png}{0.72\textwidth}
\caption{Прямая калибровка JWST F150W SBF для GO-3055 после поправки за
поглощение. Показаны четыре заранее зафиксированные выборки: все 14,
11 визуально хороших, 10 автоматически чистых и девять объектов GO-3055,
отмеченных как high-quality в Paper IV.}
\label{fig:calibration}
\end{figure*}

Для постоянной абсолютной SBF-величины ${\rm AICc}=-23.67$, для линейного
закона ${\rm AICc}=-21.22$. Постоянная модель лучше на
$\Delta{\rm AICc}=2.45$, но это не является решающим доказательством
физического отсутствия цветового члена. Leave-one-out RMS равны
соответственно $0.0955$ и $0.0989$
зв.\ вел. Линейный закон сохраняется для прямого сопоставления с Paper III,
но обнаружение цветового наклона по этой выборке не заявляется.

\begin{table*}
\centering
\caption{Линейная калибровка для четырёх заранее зафиксированных выборок.
Нуль-пункты приведены при общем $C_0=0.5675$.}
\label{tab:sample-variants}
\begin{tabular}{lrrrrr}
\toprule
Выборка & $N$ & $a$ & $b$ & $\sigma_{\rm int}$ & LOO RMS \\
\midrule
Все объекты & 14 & $-3.187$ & $0.523$ & $0.068$ & $0.099$ \\
Визуально хорошие & 11 & $-3.181$ & $0.886$ & $0.065$ & $0.105$ \\
Автоматически чистые & 10 & $-3.181$ & $0.826$ & $0.071$ & $0.130$ \\
High-quality Paper IV & 9 & $-3.161$ & $0.790$ & $0.052$ & $0.105$ \\
\bottomrule
\end{tabular}
\end{table*}

Меньший обучающий разброс подвыборки Paper IV не означает лучшей линейной
калибровки: её leave-one-out RMS больше, чем для всех 14 объектов. Для тех же
девяти точек квадратичная модель даёт RMS $0.063$ зв.\ вел., однако при девяти
объектах это чувствительный к переобучению исследовательский результат.

\subsection{Сравнение цветовых корреляций}

Для наших 14 объектов коэффициенты корреляции $\Mbar_{150,0}$ с
$(F090W-F150W)_0$ равны $r=0.22$ ($p=0.456$) и $\rho=0.05$ ($p=0.852$).
Для восьми галактик, общих с таблицей~1 Paper III, получено
$r=0.40$ ($p=0.322$). Для 14 калибраторов Paper III наблюдаемая
$\Mbar_{110}^{\rm obs}=\mbar_{110}-\mu_{\rm cluster}$ коррелирует с $g-z$
существенно сильнее: $r=0.97$ ($p<10^{-8}$), $\rho=0.96$.

При этом сами цветовые индексы для восьми общих галактик согласованы:
корреляция F090W--F150W с Pan-STARRS $g-z$ даёт
$r=0.89$ ($p=0.003$) и $\rho=0.88$ ($p=0.004$). Следовательно, слабый
SBF--цветовой тренд нельзя объяснить полной несостоятельностью инфракрасного
цвета как индикатора звёздного населения.

\begin{figure*}
\centering
\paperfigure{../../runs/sbf2_go3055/analysis/figures/go3055_color_correlation_vs_jensen.png}{0.94\textwidth}
\caption{Слева и в центре: зависимости абсолютной SBF-величины
от цвета для настоящей работы и Paper III. Справа: прямое сравнение
F090W--F150W и $g-z$ для восьми общих галактик. Наклоны первых двух панелей
нельзя сравнивать буквально из-за разных SBF-фильтров и цветовых индексов.}
\label{fig:jensen-correlation}
\end{figure*}

\subsection{Радиальные измерения}

Для внутренних колец наклон равен примерно $1.24$ при внутреннем разбросе
$0.068$ зв.\ вел.; для внешних --- $0.08$ и $0.089$ зв.\ вел.
Две точки одной галактики коррелированы и не являются 28 независимыми
калибраторами, но их различие служит диагностикой градиентов населения и
остатков моделирования.

\begin{figure*}
\centering
\paperfigure{../../runs/sbf2_go3055/analysis/figures/go3055_annulus_local_color_sbf.png}{0.78\textwidth}
\caption{Локальный цвет и абсолютная F150W SBF-величина в согласованных
кольцах. Кружки и квадраты обозначают внутреннее и внешнее кольца; серые
отрезки соединяют измерения одной галактики.}
\label{fig:annuli}
\end{figure*}

NGC~3379 имеет наиболее однородный нормированный остаток в выборке, но её
кольца различаются на $0.168$ зв.\ вел. Внешнее кольцо синее примерно на
$0.031$ зв.\ вел. и имеет более слабые SBF, причём амплитуда существенно
превышает ожидание глобального цветового закона. Поправка $P_r$ пренебрежимо
мала во внутреннем кольце и меньше одного процента $P_0$ во внешнем.
NGC~3379 сохраняется как качественное измерение и показывает, что различие
колец не равнозначно плохому остаточному изображению.

\subsection{Перекрёстная проверка по TRGB}

Для каждой галактики цветовой закон заново строится по остальным 13 объектам.
Средний остаток
\[
\langle\mu_{\rm SBF,LOO}-\mu_{\rm TRGB}\rangle=+0.0063
 \ \text{зв.\ вел.},
\]
RMS равен $0.0989$ зв.\ вел., медиана модуля остатка --- $0.0883$
зв.\ вел. Наибольшие отклонения имеют NGC~3379, NGC~4552, NGC~1399,
NGC~4472 и NGC~4649.

\begin{figure*}
\centering
\paperfigure{../../runs/sbf2_go3055/analysis/figures/go3055_sbf_vs_trgb_leave_one_out.png}{0.66\textwidth}
\caption{Модули расстояния F150W SBF, полученные с исключением проверяемой
галактики, в сравнении с TRGB тонкого среза из Paper IV. Пунктир соответствует
равенству расстояний.}
\label{fig:loo}
\end{figure*}

\subsection{Связь со светимостью галактики}

В качестве оценки светимости используется абсолютная F150W-величина гладкой
модели внутри $32.8$ угл.\ сек, а не полная величина галактики. Прямая
корреляция со SBF имеет $r=0.35$ ($p=0.220$). После вычитания линейного
цветового закона получено $r=0.48$ ($p=0.079$) и
$\rho=0.48$ ($p=0.081$). Формального обнаружения нет, но результат достаточно
близок к значимости, чтобы впоследствии заменить фиксированную апертуру
однородной полной фотометрией или оценкой звёздной массы.

\begin{figure*}
\centering
\paperfigure{../../runs/sbf2_go3055/analysis/figures/go3055_sbf_vs_host_luminosity.png}{0.92\textwidth}
\caption{Абсолютная F150W SBF-величина и остаток после цветовой поправки в
зависимости от абсолютной F150W-величины гладкой модели галактики внутри
$32.8$ угл.\ сек. Правая панель проверяет, добавляет ли светимость информацию
после учёта цвета.}
\label{fig:luminosity}
\end{figure*}

\section{Сопоставление с опубликованной калибровкой}

\begin{table}
\centering
\caption{Опубликованная и предварительная калибровки. Наклоны напрямую
несопоставимы из-за разных SBF-фильтров и цветов.}
\label{tab:numerical-comparison}
\resizebox{\columnwidth}{!}{%
\begin{tabular}{lcc}
\toprule
Величина & Paper III & Эта работа \\
\midrule
SBF-фильтр & HST F110W & JWST F150W \\
Цвет & PS $g-z$ & F090W--F150W \\
Опорный цвет & 1.30 & 0.5675 \\
Наклон & $1.86\pm0.16$ & $0.52^{+0.48}_{-0.58}$ \\
Нуль-пункт & $-2.760\pm0.016$ & $-3.187^{+0.026}_{-0.020}$ \\
Число объектов & 14 & 14 \\
LOO RMS & не приведён & 0.099 \\
\bottomrule
\end{tabular}
}
\end{table}

Геометрия измерения сознательно оставлена близкой к Paper III: используются
те же два кольца, их взвешенное среднее и цвет, соответствующий SBF-области.
Однако F150W и F090W--F150W иначе реагируют на возраст и металличность, чем
HST F110W и оптический $g-z$, поэтому настоящая работа не является буквальным
повторением калибровки.

Отличается и работа с опорными расстояниями. Средние расстояния скоплений
уменьшают шум отдельных измерений и влияние глубины скопления. Индивидуальные
расстояния Paper IV дают более строгую проверку каждой пары SBF--TRGB, но
сохраняют глубину и индивидуальный шум TRGB. До публикации необходимо показать
оба варианта, не выбирая постфактум тот, у которого меньше разброс.

\section{Обсуждение}

\subsection{Обнаружен ли цветовой член?}

Наилучший наклон полной выборки положителен, но 16-й бутстрэп-процентиль
равен $-0.054$ и пересекает ноль. Постоянная модель имеет лучший AICc и
немного меньший leave-one-out RMS. Наклон
чувствителен к балансу Девы и Форнакса, выбору кольца и узкому диапазону
цветов. Корректный вывод: линейный закон типа Paper III согласуется с
данными, но цветонезависимый нуль-пункт F150W также не исключён.

Видимая ступень Форнакс--Дева составляет около $-0.11$ зв.\ вел. Её
бутстрэп-интервал не пересекает ноль, но введение ступени ухудшает AICc
примерно на $1.05$ из-за наличия лишь трёх галактик Форнакса. Это
диагностический тест окружения, а не обнаружение второго нуль-пункта.

\subsection{Радиальное поведение и структура остатков}

Большое расхождение колец возможно при почти идеальном остаточном изображении,
а заметная крупномасштабная структура не всегда приводит к разным средним
SBF в кольцах. Поэтому контроль качества должен отдельно учитывать покрытие
моделью, структуру остатка и радиальный градиент SBF. Самая сильная
крупномасштабная структура в текущей секторной диагностике отмечена у
NGC~4472, NGC~4486, NGC~4636 и NGC~4649. У NGC~1380 и NGC~4697 неполное
покрытие внешнего кольца моделью.

\subsection{Ограничения текущего результата}

Поглощение, линеаризованный фон и общие нуль-пункты теперь явно учтены.
Остались две физически существенные группы систематики. Во-первых,
коррелированный шум i2d измерен, но ещё не включён в фит как отдельный
шаблон $N(k)$; отрицательные $P_1$ нельзя трактовать как физическую
отрицательную мощность. Во-вторых, ограничение остатка на $3.5\sigma$
изменяет медианно $3.35\%$ пикселей внутреннего кольца и даёт прокси-сдвиг
полной дисперсии $-0.092$ зв.\ вел.; во внешнем кольце соответствующие числа
равны $0.15\%$ и $-0.004$ зв.\ вел. Необходим повторный фит $P_0$ для
нескольких порогов, а не только сравнение дисперсий. Маска компактных
источников и модель полноты $P_r$ по-прежнему требуют инъекций искусственных
источников или контролируемого изменения порога.

\begin{figure*}
\centering
\paperfigure{../../runs/sbf2_go3055/analysis/figures/go3055_correlated_noise_diagnostics.png}{0.92\textwidth}
\caption{Слева: спектры мощности углов исходных F150W i2d; красная линия ---
медиана 14 полей, пунктир --- белый шум. Справа: отрицательные $P_1$ в
сопоставлении с наклоном спектра шума.}
\label{fig:correlated-noise}
\end{figure*}

\begin{figure*}
\centering
\paperfigure{../../runs/sbf2_go3055/analysis/figures/go3055_winsorization_diagnostics.png}{0.92\textwidth}
\caption{Диагностика ограничения остатка на $3.5\sigma$. Прокси-сдвиг
основан на полной дисперсии нормированного остатка и не заменяет повторный
фит спектра мощности.}
\label{fig:winsorization}
\end{figure*}

\section{Выводы}

\begin{enumerate}
\item Все 14 галактик GO-3055 прошли измерение F150W SBF и проверку
артефактов в двух кольцах Paper III. Для всех сохраняется статус ручной
проверки, главным образом из-за $P_1\leq0$.
\item Калибровка всех 14 объектов имеет нуль-пункт $-3.187$ зв.\ вел. при
$(F090W-F150W)_0=0.5675$, наклон $0.523$ и внутренний разброс $0.068$ зв.\ вел.
\item Leave-one-out проверка относительно новых расстояний Paper IV даёт
RMS $0.099$ зв.\ вел., то есть около $4.7\%$ по расстоянию, без заметного
среднего смещения.
\item Выборка пока не различает надёжно линейный цветовой закон и постоянную
абсолютную SBF-величину F150W.
\item Ни визуальный отбор 11 объектов, ни автоматический отбор 10 объектов
не улучшает leave-one-out RMS линейной модели; меньший обучающий разброс
high-quality выборки Paper IV также не переносится в более точный линейный
прогноз.
\item Внутреннее и внешнее кольца не следуют одному и тому же эмпирическому
тренду. NGC~3379 показывает, что радиальное различие SBF может сосуществовать
с исключительно чистым остатком.
\item Для прямого сопоставления с Paper III необходима параллельная
калибровка по средним расстояниям скоплений и явный учёт пересмотра
нуль-пункта TRGB в Paper IV.
\end{enumerate}

\appendix

\section{Воспроизводимость}

Измерительный ноутбук: \nolinkurl{code/sbf-2.ipynb}. Пакетный обработчик:
\nolinkurl{code/run_sbf_2_batch.py}. Итоговые таблицы и рисунки строятся в
\nolinkurl{code/sbf-2-graph.ipynb}, который читает готовые продукты и не
повторяет измерение SBF. Полная таблица расстояний в Мпк находится в
\nolinkurl{runs/sbf2_go3055/analysis/tables/go3055_distances_mpc.csv}.

\begin{thebibliography}{99}

\bibitem[Tonry \& Schneider(1988)]{TonrySchneider1988}
Tonry, J. L., \& Schneider, D. P. 1988, AJ, 96, 807.

\bibitem[Jensen et al.(2015)]{Jensen2015}
Jensen, J. B., Blakeslee, J. P., Ma, C.-P., et al. 2015,
ApJ, 808, 91, arXiv:1505.00400.

\bibitem[Jensen et al.(2021)]{Jensen2021}
Jensen, J. B., Blakeslee, J. P., Ma, C.-P., et al. 2021,
ApJS, 255, 21.

\bibitem[Jensen et al.(2025)]{Jensen2025}
Jensen, J. B., Blakeslee, J. P., Cantiello, M., et al. 2025,
\textit{The TRGB--SBF Project. III. Refining the HST Surface Brightness
Fluctuation Distance Scale Calibration with JWST},
ApJ, arXiv:2502.15935v2.

\bibitem[Chazov et al.(2026)]{Chazov2026}
Chazov, M. I., Makarov, D. I., Tully, R. B., et al. 2026,
\textit{The TRGB--SBF Project. IV. A Color Calibration of the TRGB in the
JWST F090W+F150W Filters}, arXiv:2603.11160v1.

\end{thebibliography}

\end{document}
